\section{Weighted geometric means and finite mean trees}\label{sec:matrix_mean_trees}

This section follows the first part of the transport section of
\cite{OpenAI2026AreaLaw} (Section~7, Lemma~7.1). For positive definite
matrices $A,B$ and a real parameter $p$, the weighted geometric mean is the
operation of Kubo and Ando~\cite{KuboAndo1980Means}
\begin{align}
    A\#_p B = A^{1/2}\bigl(A^{-1/2}BA^{-1/2}\bigr)^p A^{1/2}.
    \notag
\end{align}
Finite binary trees of such means carry perturbations from the leaves to the
root. Only positive definite inputs occur; the mean is never evaluated on
singular matrices.

\begin{definition}[Weighted geometric mean]\label{def:weighted_geometric_mean}
    \lean{Matrix.geomMean}
    \leanok
    For positive definite complex matrices $A,B$ and $p\in\mathbb R$, set
    $A\#_pB=A^{1/2}(A^{-1/2}BA^{-1/2})^pA^{1/2}$, where the real powers are
    taken in the functional calculus of positive definite matrices.
\end{definition}

\begin{lemma}[One-vertex properties of the mean]\label{lem:weighted_geometric_mean_vertex}
    \lean{Matrix.PosDef.geomMean, Matrix.geomMean_zero, Matrix.geomMean_one,
      Matrix.geomMean_star_conj, Matrix.geomMean_smul, Matrix.geomMean_comm,
      Matrix.geomMean_mono_right, Matrix.geomMean_mono_left,
      Matrix.re_dotProduct_rpow_mulVec_le, Matrix.re_dotProduct_geomMean_le,
      Matrix.geomMean_mul_of_commute, Matrix.commute_geomMean_of_commute}
    \leanok
    \uses{def:weighted_geometric_mean}
    Let $A,B$ be positive definite. Then $A\#_pB$ is positive definite,
    $A\#_0B=A$, $A\#_1B=B$, and:
    \begin{enumerate}
        \item $(S^*AS)\#_p(S^*BS)=S^*(A\#_pB)S$ for every invertible $S$;
        \item $(cA)\#_p(dB)=c^{1-p}d^p\,(A\#_pB)$ for $c,d>0$;
        \item $A\#_pB=B\#_{1-p}A$;
        \item for $0\le p\le1$ the mean is monotone in each argument, and for
        every vector $z$,
        $\langle z,(A\#_pB)z\rangle\le\langle z,Az\rangle^{1-p}\langle z,Bz\rangle^p$;
        \item if $A_1,B_1$ commute with $A_2,B_2$, then
        $(A_1A_2)\#_p(B_1B_2)=(A_1\#_pB_1)(A_2\#_pB_2)$, and the two factors
        commute.
    \end{enumerate}
\end{lemma}

\begin{proof}\leanok
    For an invertible $S$ put $\widehat A=S^*AS$ and
    $U=A^{1/2}S\widehat A^{-1/2}$. Then $U^*U=I$, and
    $\widehat A^{-1/2}S^*BS\widehat A^{-1/2}=U^*(A^{-1/2}BA^{-1/2})U$, so
    unitary covariance of the functional calculus gives the congruence rule.
    Scalar homogeneity follows from $(cX)^r=c^rX^r$. Congruence by
    $A^{-1/2}$ reduces exchange symmetry to the identity
    $C^p=C^{1/2}(C^{-1})^{1-p}C^{1/2}$ for $C=A^{-1/2}BA^{-1/2}$.
    Monotonicity in $B$ follows from operator monotonicity of $C\mapsto C^p$
    for $0\le p\le1$, and monotonicity in $A$ from exchange symmetry. For the
    vector inequality, write $w=A^{1/2}z$; then
    $\langle z,(A\#_pB)z\rangle=\langle w,C^pw\rangle$, and Jensen's inequality
    for the concave function $x\mapsto x^p$ against the spectral measure of $C$
    in the unit vector $w/\|w\|$ gives
    $\langle w,C^pw\rangle\le\|w\|^{2(1-p)}\langle w,Cw\rangle^p$. For the
    product rule, the inverse square root of $A_1A_2$ factors, the whitened
    second argument is the product of the two commuting positive definite
    matrices $A_g^{-1/2}B_gA_g^{-1/2}$, and the $p$th power of a product of
    commuting positive definite matrices is the product of the powers.
\end{proof}

A weighted binary tree is a finite rooted binary tree with a parameter
$p\in[0,1]$ at each internal vertex, whose two outgoing edges have weights
$1-p$ and $p$. Its leaves carry labels $j$, and the terminal weight $w_j$ of a
label is the sum, over the leaves carrying it, of the products of edge weights
along the root-to-leaf paths. Assigning positive definite matrices $A_j$ to the
labels and evaluating each internal vertex by $\#_p$ defines the root $M$. The
tree is part of the data; no associativity of matrix means is used.

\begin{definition}[Weighted binary tree of means]\label{def:mean_tree}
    \lean{Matrix.MeanTree, Matrix.MeanTree.eval, Matrix.MeanTree.weight}
    \leanok
    \uses{def:weighted_geometric_mean}
    A weighted binary tree is either a leaf with a label $j$ or an internal
    vertex with a parameter $p\in[0,1]$ and two subtrees. Its root on inputs
    $(A_j)$ is $A_j$ at a leaf and $M_l\#_pM_r$ at an internal vertex with
    subtree roots $M_l$, $M_r$. The terminal weights are $\delta_j$ at a leaf
    and $(1-p)w^{(l)}+p\,w^{(r)}$ at an internal vertex.
\end{definition}

\begin{lemma}[Terminal weights and contraction]\label{lem:mean_tree_weights}
    \lean{Matrix.MeanTree.weight_nonneg, Matrix.MeanTree.sum_weight,
      Matrix.MeanTree.eval_node_zero, Matrix.MeanTree.eval_node_one,
      Matrix.MeanTree.weight_node_zero, Matrix.MeanTree.weight_node_one}
    \leanok
    \uses{def:mean_tree}
    The terminal weights are nonnegative and sum to one. An edge of weight
    zero may be omitted and the vertex with one remaining child contracted:
    this changes neither the root nor the terminal weights.
\end{lemma}

\begin{proof}\leanok
    Induction on the tree, using $(1-p)+p=1$ and the endpoint values
    $A\#_0B=A$, $A\#_1B=B$ of Lemma~\ref{lem:weighted_geometric_mean_vertex}.
\end{proof}

\begin{theorem}[Properties of a finite mean tree]\label{thm:mean_tree_properties}
    \lean{Matrix.MeanTree.posDef_eval, Matrix.MeanTree.eval_star_conj,
      Matrix.MeanTree.eval_mono, Matrix.MeanTree.eval_smul,
      Matrix.MeanTree.re_dotProduct_eval_le, Matrix.MeanTree.smul_le_eval,
      Matrix.MeanTree.eval_listProd_ofFn, Matrix.MeanTree.commute_eval_band}
    \leanok
    \uses{def:mean_tree}
    Let a weighted binary tree have terminal weights $w_j$ and positive
    definite inputs $A_j$, and let $M$ be its root. Then $M$ is positive
    definite, commutes with invertible congruence, is monotone in the inputs,
    and satisfies
    $\mathcal T(c_jA_j:j)=\bigl(\prod_jc_j^{w_j}\bigr)M$ for $c_j>0$. For every
    vector $z$,
    \begin{align}
        \langle z,Mz\rangle\le\prod_j\langle z,A_jz\rangle^{w_j}.
        \notag
    \end{align}
    If $P$ is positive semidefinite (in particular, an orthogonal projection)
    and $A_j\ge c_jP$ with $c_j>0$, then $M\ge\bigl(\prod_jc_j^{w_j}\bigr)P$.
    Finally, if $A_j=\prod_{g=1}^KA_{j,g}$ with positive definite factors and
    every factor from band $g$ commutes with every factor from band $g'\ne g$,
    then $M=\prod_{g=1}^KM_g$, where $M_g$ is the root of the same tree with
    inputs $A_{j,g}$, and the factors $M_g$ commute.
\end{theorem}

\begin{proof}\leanok
    \uses{lem:weighted_geometric_mean_vertex, lem:mean_tree_weights}
    Positivity, congruence, monotonicity, homogeneity and the vector
    inequality follow by induction from
    Lemma~\ref{lem:weighted_geometric_mean_vertex}; in the inductive step for
    the vector inequality the exponents combine because all terminal weights
    are nonnegative. For the lower bound, fix $0<\varepsilon<1$ and choose
    $\lambda_j>0$ with $A_j\ge\lambda_jI$. Then
    $A_j-(1-\varepsilon)c_jP=(1-\varepsilon)(A_j-c_jP)+\varepsilon A_j
    \ge\varepsilon\lambda_jI$, so a common $\delta>0$ with
    $(1-\varepsilon)c_j\delta\le\varepsilon\lambda_j$ gives
    $A_j\ge(1-\varepsilon)c_j(P+\delta I)$, a positive definite input.
    Monotonicity, homogeneity, the sum rule for the terminal weights, and the
    fact that a tree with a constant input $X$ returns $X$ give
    $M\ge(1-\varepsilon)\bigl(\prod_jc_j^{w_j}\bigr)(P+\delta I)
    \ge(1-\varepsilon)\bigl(\prod_jc_j^{w_j}\bigr)P$; let $\varepsilon\to0$.
    For the band factorization, a matrix commuting with all inputs commutes
    with the root, so the roots of trees with inputs from different bands
    commute; induction over the tree with the product rule of
    Lemma~\ref{lem:weighted_geometric_mean_vertex} handles two bands, and
    induction on $K$ handles the general case.
\end{proof}

\subsection{Normalized derivative maps}

For $0<p<1$ and a positive semidefinite matrix $X$ the real power has the
L\"owner representation
$X^p=\kappa_p\int_0^\infty\bigl(t^{p-1}-t^p(t+X)^{-1}\bigr)\,dt$ with a positive
constant $\kappa_p$. A linear map on matrices is called completely positive when
its entrywise action on every matrix amplification preserves positive
semidefinite matrices.

\begin{lemma}[Derivative of a real power]\label{lem:rpow_frechet_derivative}
    \lean{Matrix.rpow_eq_integral, Matrix.hasFDerivWithinAt_rpow,
      Matrix.rpowFDeriv_apply, Matrix.isNPositiveMap_rpowFDeriv,
      Matrix.rpowFDeriv_self}
    \leanok
    Let $0<p<1$ and let $C$ be positive definite. Along Hermitian matrices, the
    map $X\mapsto X^p$ is Fr\'echet differentiable at $C$ with derivative
    \begin{align}
        D(C^p)[H]=\kappa_p\int_0^\infty t^p(t+C)^{-1}H(t+C)^{-1}\,dt.
        \notag
    \end{align}
    This derivative is completely positive and satisfies $D(C^p)[C]=pC^p$.
\end{lemma}

\begin{proof}\leanok
    Near $C$ every shift $t+X$ is invertible with
    $\|(t+X)^{-1}\|\le2(t+c)^{-1}$, where $C\ge cI$; hence the derivative of
    the integrand is dominated by a multiple of
    $t^p\bigl(t^{-1}-(t+c)^{-1}\bigr)$, which is integrable, and one may
    differentiate under the integral sign. Each integrand is a positive multiple
    of a sandwich map, so the integral is completely positive. The identity
    $D(C^p)[C]=pC^p$ is the derivative at $s=1$ of $(sC)^p=s^pC^p$.
\end{proof}

\begin{lemma}[Normalized derivative maps at a vertex]\label{lem:mean_vertex_derivative_maps}
    \lean{Matrix.hasFDerivWithinAt_geomMean_right, Matrix.hasFDerivWithinAt_geomMean_left,
      Matrix.geomMeanDerivRight_self, Matrix.geomMeanDerivLeft_self,
      Matrix.normalizedDerivMap_right_one, Matrix.normalizedDerivMap_left_one,
      Matrix.isNPositiveMap_normalizedDerivMap_right,
      Matrix.isNPositiveMap_normalizedDerivMap_left}
    \leanok
    \uses{def:weighted_geometric_mean, lem:rpow_frechet_derivative}
    Let $A,B$ be positive definite, $0\le p\le1$ and $M=A\#_pB$. At a child $D$
    of positive edge weight $w$, the map
    $\Phi_D(Z)=w^{-1}M^{-1/2}\,\partial_DM[D^{1/2}ZD^{1/2}]\,M^{-1/2}$ is unital
    and completely positive.
\end{lemma}

\begin{proof}\leanok
    \uses{lem:weighted_geometric_mean_vertex}
    By the chain rule,
    $\partial_BM[H]=A^{1/2}D(C^p)[A^{-1/2}HA^{-1/2}]A^{1/2}$ with
    $C=A^{-1/2}BA^{-1/2}$, a composition of completely positive maps; at the
    endpoints $p=0,1$ the derivative of the power is zero or the identity.
    Exchange symmetry gives $\partial_A$. The identity $D(C^p)[C]=pC^p$ gives
    $\partial_BM[B]=pM$, and symmetrically $\partial_AM[A]=(1-p)M$, which is
    unitality.
\end{proof}

\begin{theorem}[Normalized derivative maps of a mean tree]\label{thm:mean_tree_leaf_maps}
    \lean{Matrix.MeanTree.hasFDerivWithinAt_eval_update, Matrix.MeanTree.derivLabel_self,
      Matrix.MeanTree.sandwich_derivLabel_eq_smul_leafMap,
      Matrix.MeanTree.isNPositiveMap_leafMap, Matrix.MeanTree.leafMap_one,
      Matrix.MeanTree.leafMap_leaf, Matrix.MeanTree.leafMap_node_left,
      Matrix.MeanTree.leafMap_node_right}
    \leanok
    \uses{def:mean_tree, lem:mean_vertex_derivative_maps}
    Let a weighted binary tree have distinct leaves with positive definite inputs
    $A_j$, terminal weights $w_j$ and root $M$. The root is Fr\'echet
    differentiable in each input along Hermitian matrices. Composing the
    normalized vertex maps from a leaf $j$ to the root gives a completely
    positive map $\Phi_j$, the identity along an empty path, such that
    \begin{align}
        M^{-1/2}\,\partial_{A_j}M[A_j^{1/2}ZA_j^{1/2}]\,M^{-1/2}=w_j\Phi_j(Z),
        \notag
    \end{align}
    and $\Phi_j$ is unital when $w_j>0$.
\end{theorem}

\begin{proof}\leanok
    \uses{thm:mean_tree_properties}
    Since the leaves are distinct, at each vertex the input $A_j$ enters only one
    child, and the chain rule applies with the partial derivatives of
    Lemma~\ref{lem:mean_vertex_derivative_maps}. The square-root normalizations
    cancel at successive vertices, so the normalized map at the root is the
    composition of the normalized edge maps along the path. The Euler identities
    propagate to $\partial_{A_j}M[A_j]=w_jM$, which gives unitality.
\end{proof}

\begin{theorem}[Block structure of the normalized derivative maps]\label{thm:mean_tree_leaf_maps_blocks}
    \lean{Matrix.MeanTree.isBlockDiagonalMap_leafMap, Matrix.IsBlockDiagonalMap.toBlocks₁₁,
      Matrix.IsBlockDiagonalMap.toBlocks₂₂, Matrix.MeanTree.eval_fromBlocks_diag,
      Matrix.geomMean_fromBlocks_diag, Matrix.rpow_fromBlocks_diag}
    \leanok
    \uses{thm:mean_tree_leaf_maps}
    If all inputs of a weighted binary tree with distinct leaves are block
    diagonal for a fixed orthogonal decomposition, the root is block diagonal and
    the normalized leaf maps preserve each matrix block. On the first diagonal
    block they are the normalized leaf maps of the tree of first diagonal blocks,
    and they do not depend on the other diagonal blocks.
\end{theorem}

\begin{proof}\leanok
    The functional calculus of a block-diagonal Hermitian matrix acts blockwise,
    so the mean and the root are block diagonal. In the derivative integral every
    sandwich factor $(t+C)^{-1}$ is block diagonal, so the derivative of the
    power commutes with compression to each block and acts on the diagonal
    blocks as the derivatives at the diagonal blocks. These properties are
    preserved by sums, compositions, positive multiples and sandwiches by block
    diagonal matrices, which build the leaf maps.
\end{proof}

\begin{corollary}[Leaf states]\label{cor:mean_tree_leaf_states}
    \lean{Matrix.MeanTree.posSemidef_traceAdjoint_leafMap,
      Matrix.MeanTree.trace_traceAdjoint_leafMap}
    \leanok
    \uses{thm:mean_tree_leaf_maps}
    For a leaf of positive terminal weight, the trace adjoint $\Phi_j^\dagger$,
    defined by $\operatorname{Tr}(\Phi_j^\dagger(\rho)Z)=\operatorname{Tr}(\rho\,\Phi_j(Z))$,
    maps density matrices to density matrices.
\end{corollary}

\begin{proof}\leanok
    The trace adjoint of a positive map is positive, and
    $\operatorname{Tr}\Phi_j^\dagger(\rho)=\operatorname{Tr}(\rho\,\Phi_j(I))=\operatorname{Tr}\rho$
    by unitality.
\end{proof}
